%%=============================================================================
%% Resultaat en validatie
%%=============================================================================

%%TODO: toevoegen screenshots/testresults/conclusie result

\chapter{\IfLanguageName{dutch}{Resultaat en validatie}{Result and validation}}%
\label{ch:resultaat}
Dit hoofdstuk beschrijft het beoogde resultaat van het graduaatsproject en
de manier waarop de gerealiseerde proof of concept (PoC) wordt gevalideerd.
Het project richt zich op de ontwikkeling van een zelfstandige visuele
e-maileditor waarmee gebruikers de inhoud van bepaalde automatisch
gegenereerde e-mails kunnen aanpassen, opslaan, vooraf bekijken en testen.

Het beoogde resultaat is een werkende demonstratie waarin de verschillende
onderdelen van dit proces samenkomen. De gebruiker moet een bestaande
e-mailtemplate kunnen selecteren, de inhoud visueel kunnen aanpassen en
het resultaat kunnen beoordelen voordat een testmail wordt verstuurd.
Daarnaast moet de oplossing aantonen dat de gegenereerde HTML gecontroleerd
wordt verwerkt en geschikt kan worden gemaakt voor weergave in gangbare
e-mailclients.

De validatie richt zich zowel op de functionele werking als op de technische
betrouwbaarheid van de oplossing. Hierbij worden de vooraf bepaalde vereisten
vergeleken met de werkelijk gerealiseerde functionaliteiten en de resultaten
van de uitgevoerde tests.

\section{\IfLanguageName{dutch}{Wat er opgeleverd is}{What was delivered}}%
\label{sec:opgeleverd}
Het beoogde eindresultaat is een zelfstandige proof of concept van een
WYSIWYG-e-maileditor. De toepassing wordt ontwikkeld buiten de
productieomgeving van Vesalius.ai en dient om de technische haalbaarheid
van visueel aanpasbare e-mailtemplates te onderzoeken.

De PoC richt zich op drie piloottypes: afspraakbevestigingen,
afspraakherinneringen en afspraakannuleringen. Voor deze e-mails moet de
gebruiker de beschikbare inhoud en ondersteunde opmaak kunnen aanpassen
zonder rechtstreeks HTML-code te hoeven bewerken.

Naast de editor omvat het beoogde resultaat een mechanisme om templates
te laden en op te slaan. De wijzigingen moeten opnieuw beschikbaar zijn
wanneer de gebruiker een opgeslagen template opnieuw opent, binnen de
beperkingen van de gekozen opslagmethode.

Een tweede onderdeel is de renderingpipeline. Deze verwerkt de bewerkte
inhoud, controleert de HTML en verwerkt toegestane placeholders met
fictieve gegevens. Vervolgens wordt de inhoud gecombineerd met de vaste
e-mailschil, waarin onder meer het logo, de merkkleuren en de footer kunnen
zijn opgenomen. Indien nodig wordt CSS inline geplaatst om de weergave in
e-mailclients te ondersteunen.

De PoC moet daarnaast een previewfunctie aanbieden waarmee de gebruiker de
gerenderde e-mail kan controleren. Voor verdere validatie wordt een lokale
testmailomgeving voorzien, zodat de gegenereerde e-mails kunnen worden
verstuurd en geïnspecteerd zonder echte ontvangers of patiëntgegevens te
gebruiken.

Het beoogde eindresultaat bestaat bijgevolg uit de volgende onderdelen:

\begin{itemize}
    \item een visuele editor voor de inhoud en ondersteunde opmaak van
    e-mailtemplates;
    \item functionaliteit om templates te laden, aan te passen en op te slaan;
    \item een renderingpipeline voor HTML, placeholders en de vaste e-mailschil;
    \item een previewfunctie voor de gegenereerde e-mail;
    \item een lokale testmailomgeving voor het controleren van de uitvoer;
    \item geautomatiseerde tests en een evaluatie van de compatibiliteit
    met beschikbare e-mailclients;
    \item technische documentatie over de architectuur, de gemaakte keuzes
    en de beperkingen van de PoC.
\end{itemize}

In de definitieve versie van dit hoofdstuk worden screenshots opgenomen
die de werkelijk gerealiseerde functionaliteiten aantonen. Bij voorkeur
wordt de editor samen met een aangepaste template getoond, gevolgd door
een voorbeeld van de gerenderde e-mail. Indien beschikbaar kan ook een
screenshot van de ontvangen testmail in de lokale testomgeving worden
toegevoegd.

\section{\IfLanguageName{dutch}{Testen}{Testing}}%
\label{sec:testen}
De validatie moet aantonen in welke mate de PoC voldoet aan de vooraf
bepaalde functionele en technische vereisten. Hiervoor worden concrete
testscenario's opgesteld die aansluiten bij de belangrijkste gebruikersflows
en onderdelen van de renderingpipeline.

De functionele tests controleren of de gebruiker een template kan selecteren,
de inhoud kan aanpassen, wijzigingen kan opslaan en het resultaat opnieuw
kan laden. Daarnaast wordt nagegaan of de preview de verwachte inhoud
weergeeft en of de testmailfunctionaliteit correct samenwerkt met de lokale
testomgeving.

Voor de templateverwerking zijn onder meer de volgende scenario's relevant:

\begin{itemize}
    \item Een beschikbare e-mailtemplate kan worden geselecteerd en geladen.
    \item De gebruiker kan de inhoud en ondersteunde opmaak aanpassen.
    \item Wijzigingen kunnen worden opgeslagen en opnieuw worden geladen.
    \item De preview toont de verwachte gerenderde inhoud.
    \item Geldige placeholders worden vervangen door de overeenkomstige
    fictieve waarden.
    \item Onbekende of ontbrekende placeholders worden veilig afgehandeld.
    \item De vaste e-mailschil wordt correct gecombineerd met de bewerkbare
    inhoud.
    \item Een testmail kan via de lokale testomgeving worden gecontroleerd.
\end{itemize}

Naast de functionele werking wordt de veiligheid van de HTML-verwerking
gecontroleerd. Hierbij wordt nagegaan of niet-toegestane HTML-elementen
en attributen worden verwijderd of geweigerd en of dynamische waarden
veilig worden verwerkt. Ook wordt onderzocht of de renderingpipeline
onverwachte of onvolledige invoer correct afhandelt.

De compatibiliteitstests richten zich op de weergave van de gegenereerde
e-mails in Gmail, Outlook en Apple Mail, voor zover deze clients beschikbaar
zijn in de testomgeving. Hierbij worden onder meer de inhoud, de opmaak,
de vaste e-mailschil en eventuele afwijkingen gecontroleerd. Wanneer een
client niet beschikbaar is, wordt de compatibiliteit voor die client als
niet-gevalideerd vermeld.

Waar mogelijk worden de tests geautomatiseerd. Unit- en integratietests
kunnen worden gebruikt om afzonderlijke onderdelen van de verwerking te
controleren. Browsertests kunnen daarnaast de belangrijkste gebruikersflows
evalueren.

De concrete testresultaten worden geregistreerd in een afzonderlijke
testtabel. Per scenario worden de verwachte uitkomst, de werkelijke
uitkomst en de status vermeld. Hierdoor kan worden aangetoond welke
vereisten zijn geslaagd, welke problemen werden vastgesteld en welke
onderdelen bijkomende aanpassingen vereisen. De gedetailleerde resultaten
worden opgenomen in Bijlage~\ref{app:testresultaten}.

\section{\IfLanguageName{dutch}{Terugkoppeling van het bedrijf}{Feedback from the company}}%
\label{sec:terugkoppeling}
De terugkoppeling van Vesalius.ai is relevant om te beoordelen of de
oplossing aansluit bij de behoeften van het bedrijf en of de gekozen
technische aanpak bruikbaar is als basis voor verdere ontwikkeling.

Tijdens de evaluatie kan worden nagegaan of de editor voldoende
bewerkingsmogelijkheden aanbiedt, of de afbakening tussen de bewerkbare
inhoud en de vaste e-mailschil duidelijk is en of de gegenereerde e-mails
voldoen aan de verwachte functionele vereisten.

Ook de onderhoudbaarheid van de oplossing is van belang. De scheiding tussen
de frontend, de mock-API, de opslag en de renderingpipeline moet het mogelijk
maken om de afzonderlijke onderdelen te begrijpen en verder te ontwikkelen.
Daarnaast moet worden beoordeeld of de gemaakte keuzes aansluiten bij de
bestaande technische context van Vesalius.ai.

De concrete feedback van de bedrijfsmentor en de betrokken teamleden wordt
in de definitieve versie van dit hoofdstuk opgenomen zodra deze beschikbaar
is. Daarbij worden de belangrijkste opmerkingen, eventuele voorgestelde
aanpassingen en de mogelijke vervolgstappen beschreven.

De PoC is in de eerste plaats bedoeld als technische demonstratie. Een
positieve evaluatie betekent daarom niet automatisch dat de oplossing
klaar is voor productiegebruik. Daarvoor is een afzonderlijke beoordeling
van de integratie, beveiliging, opslag, het beheer en de operationele
vereisten noodzakelijk.

\section{\IfLanguageName{dutch}{Beperkingen}{Limitations}}%
\label{sec:beperkingen}
De belangrijkste beperking van het project is de afgebakende scope. De PoC
richt zich op drie specifieke e-mailtypes en op het aanpassen van de
e-mailinhoud. Het volledig vrij ontwerpen van e-maillay-outs valt buiten
de doelstelling. De vaste e-mailschil blijft grotendeels behouden.

Daarnaast wordt de toepassing niet geïntegreerd in de productieomgeving
van Vesalius.ai. De mock-API en de gekozen opslagmethode dienen om de
functionaliteiten in een gecontroleerde testomgeving te demonstreren.
Ze vormen niet noodzakelijk een geschikte oplossing voor productiegebruik.

Ook de testgegevens zijn bewust beperkt tot fictieve informatie. Hierdoor
kunnen de belangrijkste verwerkingsstappen worden onderzocht zonder echte
patiëntgegevens te gebruiken. De PoC toont daarmee niet automatisch aan
dat alle vereisten voor een productieomgeving zijn ingevuld.

Een verdere beperking betreft de compatibiliteit met e-mailclients.
HTML en CSS kunnen verschillend worden geïnterpreteerd door Gmail, Outlook
en Apple Mail. Zelfs wanneer de preview en de uitgevoerde tests correct
verlopen, kan niet zonder bijkomende validatie worden gegarandeerd dat
iedere e-mailclient de inhoud identiek weergeeft.

Ten slotte hangt de geldigheid van de conclusies af van de werkelijk
uitgevoerde tests en de beschikbare testomgeving. Niet-uitgevoerde
scenario's en niet-beschikbare e-mailclients worden daarom expliciet als
beperkingen vermeld.

De resultaten moeten bijgevolg worden geïnterpreteerd als een evaluatie van
de gerealiseerde PoC binnen de onderzochte context. Ze vormen een basis
om te bepalen welke onderdelen voldoende zijn aangetoond en welke verdere
ontwikkeling of validatie nodig hebben voordat een mogelijke integratie
in Vesalius.ai kan worden overwogen.
